\documentclass{article}

\usepackage{hyperref}
\usepackage{listings}
\usepackage[dvipsnames]{xcolor}
\usepackage{amsmath}

\title{}
\author{Ryan Baker}
\date{\today}

\definecolor{JadeGreen}{RGB}{91,158,62}
\lstdefinestyle{catppuccin}{
	backgroundcolor=\color{White},
	commentstyle=\color{Gray},
	numberstyle=\footnotesize\ttfamily\color{Gray},
	stringstyle=\color{JadeGreen},
	keywordstyle=\color{BurntOrange},
	basicstyle=\ttfamily\footnotesize\color{Black},
	breakatwhitespace=false,
	breaklines=true,
	captionpos=b,
	keepspaces=true,
	numbers=left,
	numbersep=5pt,
	showspaces=false,
	showstringspaces=false,
	showtabs=false,
	tabsize=4,
}
\lstset{style=catppuccin}

\begin{document}

\maketitle
\tableofcontents
\pagebreak

\subsection*{Lecture Objectives}

\section{Functions}

\begin{itemize}
	\item What is a function?
	\begin{itemize}
		\item In programming, a function is a reusable block of code
		\item It (optionally) takes input and (optionally) returns an output
	\end{itemize}
	\item Why use functions?
	\begin{itemize}
		\item We often want to repeat the same behavior on different pieces of data
		\item Rather than pasting the same code many times, we use a function
		\begin{itemize}
			\item Functions help to keep code maintainable and readable
		\end{itemize}
		\item There is a balance to strike when extracting code into functions
		\begin{itemize}
			\item Too few functions results in long and repetitive code
			\item Too many functions will result in sub-optimal performance and a code base that is very hard to read \begin{itemize}
				\item Every time a function is called a new frame needs to be pushed to the stack and we jump around the executable
			\end{itemize}
		\end{itemize}
	\end{itemize}
	\item How to define a function: \textbf{type name(arguments)}
	\begin{itemize}
		\item \textbf{type}: The return type of the function (can be \texttt{void})
		\item \textbf{name}: The function's name
		\item \textbf{arguments}: The input arguments to a function
		\begin{itemize}
			\item Specified as \textbf{type name} in a comma seperated list
			\item \textbf{Example}: \texttt{int add(int a, int b, int c) \{...\}}
		\end{itemize}
		\item Together, the function's name and arguments make up the signature
	\end{itemize}
	\item How to call a function: \texttt{name(arguments)}
	\begin{itemize}
	\item \textbf{Example}: \texttt{int sum = add(1, 2, 3); // sum = 6}
	\end{itemize}
	\lstinputlisting[language=C++]{function.cpp}
\end{itemize}

\section{Conditions and Branches}

\noindent
Often times, we want to evaluate a certain condition, and, depending on the result, decide which code to execute.

\begin{itemize}
	\item What is branching?
	\begin{itemize}
		\item Branching is when we jump, or ``branch'' to another part of our code depending on the result of some conditional
		\item The computer has to load a new part of the executable to memory
	\end{itemize}
	\item Conditional statements
	\begin{itemize}
		\item Conditional statements are expressions that return a \textit{boolean} value
		\item Often in C++, we use comparison operators to form conditionals
		\begin{itemize}
			\item \texttt{==}: returns true if the arguments are equal (\texttt{a == b})
			\item \texttt{!=}: returns true if the arguments aren't equal (\texttt{a != b})
			\item \texttt{<}: true if left operand is less than right operand (\texttt{a < b})
			\item \texttt{>}: true if left operand is greater than right operand	 (\texttt{a > b})
			\item \texttt{<=}: true if left operand is less than or equal to right operand
			\item \texttt{>}: true if left operand is greater than or equal to right operand	
			\item \texttt{bool(x)}: returns true if \texttt{x != 0}
		\end{itemize}
		\lstinputlisting[language=C++]{conditional.cpp}
	\end{itemize}
\end{itemize}

\subsection{\texttt{if} statements}

\begin{itemize}
	\item \texttt{if} statement syntax: \texttt{if (condition) \{...\}}
	\begin{itemize}
		\item Only executes the code within \texttt{\{...\}} if the condition is true
	\end{itemize}
	\item \texttt{else} statements: \texttt{if (condition) \{...\} else \{...\}}
	\begin{itemize}
		\item Executes the first block if the condition is true, else the second block
	\end{itemize}
	\item \texttt{\{\}} is not strictly required if the block is one line
	\lstinputlisting[language=C++]{ifelse.cpp}
	\item Chaining \texttt{if} and \texttt{else} statements:
	\lstinputlisting[language=C++]{elseif.cpp}
\end{itemize}

\subsection{\texttt{switch} statements}

\begin{itemize}
	\item Switch statements jump to a case depending on the result of the expression
	\item \texttt{switch} statement syntax: \texttt{switch (expression) \{ cases... \}}
	\lstinputlisting[language=c++]{switch.cpp}
	\item If no \texttt{case} evaluates to true, the \texttt{default} case is executed
	\item Without \texttt{break} statements, every case below will also execute
\end{itemize}

\section{Loops}

\subsection{\texttt{for} loops}

\subsection{\texttt{while} loops}

\end{document}